% Part: first-order-logic
% Chapter: tableaux
% Section: provability-quantifiers

\documentclass[../../../include/open-logic-section]{subfiles}

\begin{document}

\olfileid[fr]{fol}{tab}{qpr}
\olsection{\usetoken{S}{derivability} et quantificateurs}

\begin{explain}
La démonstration du théorème de complétude exige aussi que les règles
des tableaux donnent les propriétés de~$\Proves$ établies dans cette
section.
\end{explain}

\begin{thm}
\ollabel{thm:strong-generalization} Si $c$ est une constante qui ne
figure ni dans $\Gamma$ ni dans~$!A(x)$ et si $\Gamma \Proves !A(c)$,
alors $\Gamma \Proves \lforall[x][!A(x)]$.
\end{thm}

\begin{proof}
Supposons que $\Gamma \Proves !A(c)$. Il existe donc
$!B_1, \dots, !B_n \in \Gamma$ et un !!{tableau} fermé pour
\begin{align*}
\{ \sFmla{\False}{!A(c)}, &
\sFmla{\True}{!B_1}, \dots, \sFmla{\True}{!B_n} \}.
\intertext{Il faut montrer qu'il existe aussi un !!{tableau} fermé pour}
\{ \sFmla{\False}{\lforall[x][!A(x)]}, &
\sFmla{\True}{!B_1}, \dots, \sFmla{\True}{!B_n} \}.
\end{align*}
Dans le !!{tableau} fermé donné, remplaçons la première hypothèse par
$\sFmla{\False}{\lforall[x][!A(x)]}$, puis insérons
$\sFmla{\False}{!A(c)}$ après les hypothèses.
\begin{center}
\begin{tableau}{not line numbering}
  [\sFmla{\False}{\formula{A}(c)}
    [\sFmla{\True}{\formula{B}_1}
      [\vdots
      [\sFmla{\True}{\formula{B}_n}, 
        [][][]]]]]
\end{tableau}\qquad
\begin{tableau}{not line numbering}
  [\sFmla{\False}{\lforall[x][\formula{A}(x)]}
    [\sFmla{\True}{\formula{B}_1}
      [\vdots
        [\sFmla{\True}{\formula{B}_n},
          [\sFmla{\False}{\formula{A}(c)}
        [][][]]]]]]
\end{tableau}
\end{center}
Le tableau demeure fermé, puisque tous les !!{sentence}s qui figuraient
auparavant parmi les hypothèses sont encore disponibles en haut du
!!{tableau}. La ligne insérée résulte d'une application correcte de
$\TRule{\False}{\lforall}$ : la !!{constant}~$c$ ne figure ni dans
$!B_1$, \dots, $!B_n$, ni dans $\lforall[x][!A(x)]$, donc elle ne
figure nulle part au-dessus de la ligne insérée dans le nouveau
!!{tableau}.
\end{proof}

\begin{prop}
\ollabel{prop:provability-quantifiers}
Pour tout terme clos~$t$, on a :
\begin{tagenumerate}{prvEx,prvAll}
\tagitem{prvEx}{$!A(t) \Proves \lexists[x][!A(x)]$.}{}
\tagitem{prvAll}{$\lforall[x][!A(x)] \Proves !A(t)$.}{}
\end{tagenumerate}
\end{prop}

\begin{proof}
\begin{tagenumerate}{prvEx,prvAll}
\tagitem{prvEx}{Voici un !!{tableau} fermé pour
  $\sFmla{\False}{\lexists[x][!A(x)]}, \sFmla{\True}{!A(t)}$ :
  \begin{oltableau}
    [\sFmla{\False}{\lexists[x][\formula{A}(x)]}, just = \TAss
      [\sFmla{\True}{\formula{A}(t)}, just = \TAss
        [\sFmla{\False}{\formula{A}(t)}, just = {\TRule{\False}{\lexists}[1]}, close]]]
    \end{oltableau}}{}
\tagitem{prvAll}{Voici un !!{tableau} fermé pour
  $\sFmla{\False}{\formula{A}(t)},
  \sFmla{\True}{\lforall[x][!A(x)]}$ :
  \begin{oltableau}
    [\sFmla{\False}{\formula{A}(t)}, just = \TAss
      [\sFmla{\True}{\lforall[x][\formula{A}(x)]}, just = \TAss
        [\sFmla{\True}{\formula{A}(t)}, just = {\TRule{\True}{\lforall}[2]}, close]]]
    \end{oltableau}}{}
\end{tagenumerate}
\end{proof}

\begin{editorial}
La fermeture du terme~$t$, exigée par les règles de quantification de
ce système, est rappelée explicitement dans l'énoncé de la proposition.
\end{editorial}

\end{document}
